\section*{第 \ref{ch:g10:why-clouds-form} 章——云为什么会形成}

\begin{solution}{exo:g10:why-clouds-form:1}
\omterm{def:g10:why-clouds-form:humidity}{相对湿度}是空气中水汽含量与
同温同压下\omterm{def:g10:why-clouds-form:humidity}{饱和}含水量之比，以百分数表示。
\omterm{def:g10:why-clouds-form:dew-point}{露点}是空气在保持含水量和气压不变的条件下，
冷却到\omterm{def:g10:why-clouds-form:humidity}{饱和}时所达到的温度。
\end{solution}

\begin{solution}{exo:g10:why-clouds-form:2}
\num{3.8}、\num{7.6}、\num{14.7} 和 \qty{27.2}{g/kg}。每升高
\qty{10}{\degreeCelsius}，数值约增至两倍（$7.6/3.8 = 2.0$，
$14.7/7.6 \approx 1.9$，$27.2/14.7 \approx 1.9$）。
\end{solution}

\begin{solution}{exo:g10:why-clouds-form:3}
在 \qty{24}{\degreeCelsius} 时的\omterm{def:g10:why-clouds-form:humidity}{饱和}含水量为 \qty{18.9}{g/kg}，\omterm{def:g10:why-clouds-form:humidity}{相对湿度}
为 $10.0/18.9 \approx \qty{53}{\percent}$。
\omterm{def:g10:why-clouds-form:humidity}{饱和}含水量为 \qty{10.0}{g/kg} 时的温度为 \qty{14}{\degreeCelsius}，即\omterm{def:g10:why-clouds-form:dew-point}{露点}。
\end{solution}

\begin{solution}{exo:g10:why-clouds-form:4}
夏季空气潮湿，其\omterm{def:g10:why-clouds-form:dew-point}{露点}远高于玻璃杯的温度，
与玻璃接触的空气被冷却到\omterm{def:g10:why-clouds-form:dew-point}{露点}以下，水汽在杯壁上凝结。
室内冬季空气非常干燥，
其\omterm{def:g10:why-clouds-form:dew-point}{露点}低于玻璃杯的温度：与之接触并被冷却后，
空气仍未达到\omterm{def:g10:why-clouds-form:humidity}{饱和}。
\end{solution}

\begin{solution}{exo:g10:why-clouds-form:5}
从下方加热（晴朗地面上午后的积云）；山脉（
气流被抬升越过山脉）；锋面（暖空气爬升至冷空气之上）；辐合
（气流在赤道附近汇聚，或流入低压区）。
\end{solution}

\begin{solution}{exo:g10:why-clouds-form:6}
$125 \times (30 - 18) = \qty{1500}{m}$; $125 \times (22 - 20) =
\qty{250}{m}$。第二天：云仅在距地面 \qty{250}{m} 处
形成。
\end{solution}

\begin{solution}{exo:g10:why-clouds-form:7}
在 \qty{1}{km} 处：$18 - 9.8 = \qty{8.2}{\degreeCelsius}$；在 \qty{2}{km} 处：
$18 - 19.6 = \qty{-1.6}{\degreeCelsius}$。下沉时，空气以相同速率增温，
并在地面再次达到 \qty{18}{\degreeCelsius}。
\end{solution}

\begin{solution}{exo:g10:why-clouds-form:8}
当递减率为 \qty{4}{\degreeCelsius/km} 时，未\omterm{def:g10:why-clouds-form:humidity}{饱和}气块
（\qty{9.8}{\degreeCelsius/km}）与\omterm{def:g10:why-clouds-form:humidity}{饱和}气块
（\qty{5}{\degreeCelsius/km}）降温均快于周围环境，两种情况下均稳定。
当递减率为 \qty{12}{\degreeCelsius/km} 时，即使未\omterm{def:g10:why-clouds-form:humidity}{饱和}的气泡，
降温也比周围环境慢，它保持较高温度并继续上升，
暖空气泡由此自由上升。
\end{solution}

\begin{solution}{exo:g10:why-clouds-form:9}
$14.7 - 7.6 = \qty{7.1}{g}$（每千克空气）。水汽凝结时释放的\omterm{def:g10:earth-energy-budget:latent-heat}{潜热}
使凝结处的空气升温，减缓其降温。
\end{solution}

\begin{solution}{exo:g10:why-clouds-form:10}
体积与直径的三次方成正比：$(0.5/0.02)^3 = 25^3 = 15\,625$
个小液滴对应毛毛雨滴；$(2/0.02)^3 = 100^3 = 1\,000\,000$ 对应
雨滴。
\end{solution}

\begin{solution}{exo:g10:why-clouds-form:11}
晴朗的夜空下，地面向外辐射红外线而降温，与之接触的空气也随之降温。
这种冷空气比上方空气密度大，沿斜坡向下流动，
聚集在谷底。在谷底，空气整夜持续降温，直至达到\omterm{def:g10:why-clouds-form:dew-point}{露点}，
水汽在空气中凝结成小液滴，即形成雾。
山顶上的空气没有那么冷，
仍高于\omterm{def:g10:why-clouds-form:dew-point}{露点}。
\end{solution}

\begin{solution}{exo:g10:why-clouds-form:12}
设高度为 $z$ 处的温度为 $T - a z$，\omterm{def:g10:why-clouds-form:dew-point}{露点}为 $T_d - b z$；
二者相等时，$(a - b) z = T - T_d$，故 $z = (T - T_d)/(a - b)$。取
$b = 1.7$，则 $a - b = 8.1$，每度对应 $1000/8.1 \approx \qty{123}{m}$。
\end{solution}

\begin{solution}{exo:g10:why-clouds-form:13}
\omterm{def:g10:why-clouds-form:condensation-level}{凝结高度}：$125 \times (15 - 11) = \qty{500}{m}$；该处温度：
$15 - 9.8 \times 0.5 \approx \qty{10.1}{\degreeCelsius}$。山顶：$10.1 - 5
\times 1.5 = \qty{2.6}{\degreeCelsius}$。以\omterm{def:g10:why-clouds-form:adiabatic}{干绝热递减率}下降 \qty{2}{km}：
$2.6 + 9.8 \times 2 \approx \qty{22.2}{\degreeCelsius}$，
比迎风一侧约高七度。
\end{solution}

\begin{solution}{exo:g10:why-clouds-form:14}
在 \qty{-20}{\degreeCelsius} 时：$126/103 \approx 1.22$，
对冰过\omterm{def:g10:why-clouds-form:humidity}{饱和} \qty{22}{\percent}；在 \qty{-10}{\degreeCelsius} 时：$287/260 \approx
1.10$，过\omterm{def:g10:why-clouds-form:humidity}{饱和} \qty{10}{\percent}。在这些温度中，云越冷，
冰的优势越大：该过程在云的上部寒冷区域效果最好，
那里小液滴与冰晶混合存在。
\end{solution}

\begin{solution}{exo:g10:why-clouds-form:15}
未\omterm{def:g10:why-clouds-form:humidity}{饱和}时：气块降温 \qty{9.8}{\degreeCelsius}，周围环境降温
\qty{6.5}{\degreeCelsius}，气块比周围冷 \qty{3.3}{\degreeCelsius}，
故下沉。\omterm{def:g10:why-clouds-form:humidity}{饱和}时：气块降温 \qty{5}{\degreeCelsius}，
因而比周围暖 \qty{1.5}{\degreeCelsius}，继续上升。空气只有在一个条件下才不稳定，
即上升的空气已经\omterm{def:g10:why-clouds-form:humidity}{饱和}。
\end{solution}

\begin{solution}{pb:g10:why-clouds-form:1}
\textbf{1.} \qty{14.7}{g/kg}。

\textbf{2.} $11.4/14.7 \approx \qty{78}{\percent}$。

\textbf{3.} \qty{16}{\degreeCelsius}，此时\omterm{def:g10:why-clouds-form:humidity}{饱和}含水量为
\qty{11.4}{g/kg}。

\textbf{4.} $11.4/24.1 \approx \qty{47}{\percent}$。

\textbf{5.} 不对：其含水量仍为 \qty{11.4}{g/kg}。\omterm{def:g10:why-clouds-form:humidity}{相对湿度}
下降，是因为气温升高后，它能容纳更多水汽。

\textbf{6.} 在其\omterm{def:g10:why-clouds-form:dew-point}{露点}，即 \qty{16}{\degreeCelsius}。

\textbf{7.} $11.4 - 8.7 = \qty{2.7}{g}$，即每千克空气。

\textbf{8.} 晴朗的夜空下，草地和屋顶向外辐射红外线，
温度低于周围空气；与它们接触的空气
先于其余空气降至\omterm{def:g10:why-clouds-form:dew-point}{露点}以下。

\textbf{9.} $(28.0 - 13.4)/1.5 \approx \qty{9.7}{\degreeCelsius/km}$，
接近\omterm{def:g10:why-clouds-form:adiabatic}{干绝热递减率}：午后的对流泡已使该层混合。

\textbf{10.} $(13.4 - 0.4)/2.0 = \qty{6.5}{\degreeCelsius/km}$。

\textbf{11.} 对未\omterm{def:g10:why-clouds-form:humidity}{饱和}气块而言是稳定的，因为其降温较快
（\qty{9.8}{\degreeCelsius/km}）；对\omterm{def:g10:why-clouds-form:humidity}{饱和}气块而言是不稳定的，因为其
降温较慢（\qty{5}{\degreeCelsius/km}）。

\textbf{12.} 气温随高度升高，从 \num{-0.3} 升至
\qty{1.2}{\degreeCelsius}：这是一个非常稳定的层，起到盖子的作用。
上升的云在此停止发展，云顶高度约为 \qty{4}{km}。

\textbf{13.} $28 - 16 = \qty{12}{\degreeCelsius}$。

\textbf{14.} 气温 $28 - 9.8 = \qty{18.2}{\degreeCelsius}$；\omterm{def:g10:why-clouds-form:dew-point}{露点}
$16 - 1.8 = \qty{14.2}{\degreeCelsius}$。

\textbf{15.} 由 12 降至 $18.2 - 14.2 = 4.0$：差值缩小了
\qty{8}{\degreeCelsius}，速率为每千米 \qty{8}{\degreeCelsius}。

\textbf{16.} 每相差一度，需在 $1/8$ km 内消除，即
\qty{125}{m}。

\textbf{17.} $125 \times (20 - 16) = \qty{500}{m}$。

\textbf{18.} \omterm{def:g10:why-clouds-form:dew-point}{露点}保持在 \qty{16}{\degreeCelsius}，而近地面空气
升温，温差增大，
消除温差的高度也随之升高。

\textbf{19.} $28 - 9.8\,z = 16 - 1.8\,z$。

\textbf{20.} $8\,z = 12$，故 $z = 1.5$：午后积云的底部
位于平原以上\textbf{约 \qty{1.5}{km}} 处（用水的性质进行精确计算，得
\qty{1.51}{km}）。
\end{solution}
